%! End Behavior and Asymptotes — Unit 2, Lesson 2.5 — Warm-Up
\documentclass[10pt]{article}
\usepackage{atda-article}
\usepackage{atda-boxes}
\newcommand{\TallMath}[1]{$\displaystyle #1\rule[-1.4em]{0pt}{3.2em}$}
\setlength{\workrowsep}{6pt}   % handwriting room; moves blank and key together

\begin{document}

\pageheader{Unit 2, Lesson 2.5}{Warm-Up}
% NAMESTRIP: no \namedateperiod here. The student writes their name once, on the
% cover this component is stapled behind. See references/conventions.md.

\vspace{0.15in}

\begin{notesbox}{Halving Forever, and a Number That Is Never Negative}
\small
Five quick items. Items 1 and 4 both ask a question you will answer properly today: \emph{can this
ever get to zero?}

\begin{enumerate}[leftmargin=1.6em, itemsep=4pt, topsep=4pt]

\item \textbf{Keep halving.} Fill in the next three numbers.
\smallskip

\renewcommand{\arraystretch}{1.3}
\begin{tabular}{@{}|c|c|c|c|c|c|c|@{}}
\hline
Step & $0$ & $1$ & $2$ & $3$ & $4$ & $5$ \\ \hline
Value & $64$ & $32$ & $16$ & $8$ & \blank{1.0cm} & \blank{1.0cm} \\
\hline
\end{tabular}
\smallskip

(a) Give the value at step $6$: \blank{2.0cm} \quad (b) Will this list ever reach exactly $0$?
Answer in one sentence.

\writelines{1}

\item \textbf{Last night's homework.} Item 10 graphed coffee cooling in an insulated flask. What
number did the temperatures get closer and closer to, and did the graph ever reach it?

\writelines{1}

\item \textbf{An input the rule will not take} (Lessons 1.7 and 2.1). For which input is
$\dfrac{4}{x-2}$ undefined? Then write the domain of that expression in words.

\writelines{1}

\item \textbf{Evaluating an exponential} (Lesson 2.2). Let $g(x)=2^x$. Fill in every cell.
\smallskip

\renewcommand{\arraystretch}{1.3}
\begin{tabular}{@{}|c|c|c|c|c|@{}}
\hline
$x$ & $3$ & $0$ & $-2$ & $-4$ \\ \hline
$g(x)=2^x$ & \blank{1.2cm} & \blank{1.2cm} & \blank{1.2cm} & \blank{1.2cm} \\
\hline
\end{tabular}
\smallskip

Is any one of those four outputs zero or negative? \blank{2.2cm}

\item \textbf{Infinity in interval notation} (Lesson 2.1). \quad (a) Write ``all real numbers greater
than $20$'' in interval notation. \quad (b) Write $(-\infty,0)$ as an inequality.

\writelines{1}

\end{enumerate}
\end{notesbox}

\end{document}
